\subsection{Applications of the common support to the actual RH points}
\label{subsec:rh-four-point-applications}

We use the actual RH datum of
Theorem~\ref{thm:rh-symplectic-arithmetic-data} with $e\geq2$.
The following argument supplies the three point-arithmetic
premises of the integral comparison, using the preceding
reconstruction of the ordinary-family, seed, and component
suppliers. It only requires a good-prime purity window, not
the full local statement \texttt{WitA} of the original source.

\begin{lemma}[A nonsplit symplectic extension of a block endpoint]
\label{lem:rh-four-block-extension}
Let an original rational \texttt{KPair} have a block-diagonal
representation $r$ for the blocks $U=\{0,2\}$ and $V=\{1,3\}$.
Suppose its first block is absolutely irreducible. In either
coefficient characteristic there exists an original \texttt{KPair}
with the same characteristic polynomials at every group element
whose representation satisfies \texttt{BlockExt}, hence
\texttt{Star}.
\end{lemma}

\begin{proof}
Work over a complete DVR $D$ with finite residue field
$k_D$ of order $q$, uniformizer $\pi$, and the original
integral representation and flag; a finite enlargement
is harmless. The \texttt{Small} condition and the common
coefficient construction above allow this, with all original
characters realized in the same ring. In particular in
mixed characteristic the given local embedding of the
coefficient DVR into the ambient field is compatible with
the embedding of $\mathbb Z_\ell$: the latter is uniquely
determined by density of $\mathbb Z$ in its completion
and continuity of the local embeddings into the rank-one
valued algebraic extension of the complete fraction field.
In equal characteristic one uses the unique finite
constant-field lifts. We write $J_{\rm sp}$ for the
original alternating form.

The blocks are complementary Lagrangians, and $V$ carries
the contragredient representation of $U$ twisted by the
similitude. Thus both blocks are absolutely irreducible.
In the symplectic Lie algebra let
\[
 M=\{Z:Z|_V=0,\quad Z(U)\subset V,\quad
                Z^\dagger=-Z\}.
\]
It is the rank-three free module of symmetric maps
from $U$ to its paired dual, is stable for
$\operatorname{Ad}r$, and $ZZ'=0$ for $Z,Z'\in M$.
The residual characters on this module are
$\bar\alpha^{-1},\bar\beta,\bar\tau^{-1}$, since its
independent entries are $10,12,30$ in the original
ordering. Complex conjugation $c$ acts on $U$ as $-1$
and on $V$ as $+1$. Indeed its square is one,
$2$ is invertible, and the unwanted eigensummands
have zero reduction. Consequently $\operatorname{Ad}r(c)$
acts on $M$ by $-1$.

The ordinary flag splits integrally along the blocks:
the unit resultant of the two block characteristic
polynomials at the original local element $\mathrm{ph}$
gives the two complementary projectors as polynomials
in $r(\mathrm{ph})$. These projectors preserve every
member of the flag. All its successive quotients
are free, and (BU8) separates the two residual
blocks at this element. The graded lines thus belong
to $U,V,U,V$ in that order. Equivalently, after choosing
bases of each block adapted to the split flag, the
full flag has the standard interleaved ordering.
Let $P$ be the submodule of $M$ preserving this
flag. In the adapted bases the only independent
permitted entry is $12$; the symmetric-map condition
sets the other two independent entries to zero.
Hence $P$ is a rank-one direct summand and $M/P$
is free of rank two. Both are invariant under the
local decomposition group.

Write $M_j=M/\pi^jM$ and $P_j=P/\pi^jP$. All
these are finite $\ell$-primary Galois modules.
The original local element $s_\ell$ of (BU8) has
cyclotomic value one and none of the three residual
characters of $M$ is one there. In particular
the endomorphism $s_\ell-1$ is invertible on each
of $M_j$ and $M_j/P_j$. The same is true on their
Tate duals, since the cyclotomic value is one.
This proves the necessary local $H^0$ vanishings
and, by local Tate duality, the local $H^2$
vanishings at $\ell$. It also proves the global
$H^0$ vanishing for the Tate dual of $M_j$.
The global $H^0(M_j)$ is zero by the action of $c$.

At each $p\in S\setminus\{\ell\}$ put
$N_p=I_p\cap\ker\bar r_0$. By the original killing
condition (here $\mathrm{Sb}=\varnothing$), $r$ kills
$N_p$. The local quotient by $N_p$ is an extension
of the procyclic unramified group by finite
prime-to-$\ell$ inertia. Its $\ell$-cohomological
dimension is at most one, and its $H^1$ on $M_j$
has the same cardinality as $H^0$ (take inertia
invariants and apply the procyclic calculation).
Use its inflated $H^1$ as the permitted local
subgroup. Outside $S$ impose the unramified
condition. At $\ell$ take the kernel of
\[
 H^1(\mathbb Q_\ell,M_j)\longrightarrow
 H^1(\mathbb Q_\ell,M_j/P_j).
\]
The local Euler characteristic formula, with the
vanishings just checked, gives cardinalities $q^{3j}$
and $q^{2j}$ respectively. Thus the kernel has
cardinality at least $q^j$. At the real place the
cohomology in degree one is zero and the invariants
are zero, since the action is $-1$ and $\ell$ is odd.

The finite Poitou--Tate Selmer cardinality formula
now gives at least $q^j$ global classes with these
local conditions. Indeed the ratio of the Selmer
group order to the order of the dual Selmer group
is the global invariant ratio times the product
of the local ratios (permitted subgroup order
divided by invariant order). The global ratio
is one; the local ratios away from $\ell$ and
at infinity are one; at $\ell$ the ratio is at
least $q^j$. No characteristic-zero
automorphy or finiteness of a Selmer lattice has
been used in this estimate.

Normalize each such cocycle $z_j$ by $z_j(c)=0$.
There is exactly one representative of the class
with this normalization because $1-\operatorname{Ad}r(c)=2$
on $M_j$. The condition at $\ell$ is equivalent
to the existence of a unique $\bar y_j\in M_j/P_j$
with
\[
 z_j(s)=\operatorname{Ad}r(s)\bar y_j-\bar y_j
                       \pmod {P_j}\quad(s\in G_{\mathbb Q_\ell}).
\]
Uniqueness follows from the local invariant vanishing.
The inertia conditions away from $\ell$ mean that
$z_j$ itself kills the prescribed inertia kernels,
since those kernels act trivially. Therefore the
set $E_j$ of these normalized pairs is finite,
has at least $q^j$ elements, and is stable under
reduction from the higher levels. Call a pair
primitive if its reduction at level one is nonzero.

There is a primitive pair at every level. To verify
this carefully, multiplication by $\pi^{m-n}$
identifies level-$n$ pairs with level-$m$ pairs
divisible by $\pi^{m-n}$. Divisibility of a pair
means divisibility both of all the cocycle values
and of $\bar y$; division respects all the displayed
equations because the modules $M,M/P$ are free.
If there were no primitive pair at level $n$, any
pair at level $m\ge n$ would be divisible by $\pi$.
Repeatedly divide and reduce to level $n$ until
all level-$m$ pairs have been identified with
pairs of level at most $n-1$. This would bound
$|E_m|$ by $|E_{n-1}|$ for every $m$, contradicting
$|E_m|\ge q^m$ (put $E_0=\{0\}$). The finite sets
of primitive pairs at each level now have a
compatible inverse-limit element: the finite
intersection property follows by restricting
one pair from any sufficiently high level.
This produces a continuous $M$-valued cocycle
$z$ and $\bar y\in M/P$. They satisfy the integral
conditions, and $z\ne0$ by the uniqueness of
the local coboundary parameter. Choose an integral
lift $y$ of $\bar y$. Multiply $z,y$ by $\pi$,
so that the residual representation and flag will
remain unchanged.

Define the new representation and ordinary flag
matrix by
\[
 r'(g)=(1+z(g))r(g),\qquad u'=(1-y)u .
\]
Because $M^2=0$ and $M$ is adjoint-stable,
the cocycle law is precisely the representation
law for $r'$. The factor $1+z(g)$ is symplectic:
the quadratic term in its symplectic identity
vanishes by $M^2=0$, and the linear term by
$z^\dagger=-z$. The multiplier is unchanged.
Similarly $1-y$ is symplectic. Conjugating at
$\ell$ by $1-y$ replaces the extension term by
$z(s)+y-\operatorname{Ad}r(s)y\in P$; hence the
original ordered flag is stable and the graded
characters are unchanged (the action of $P$ on
its graded lines is zero). The required inertia
kernels are still killed, and reduction modulo
$\pi$ leaves the original diagonal and flag
characters. The new pair is \texttt{Small}, with
the same complete DVR. All \texttt{KPair} conditions
are retained.

The extension has the same characteristic
polynomials, by its block triangular shape.
If it split by conjugation with a block off-diagonal
matrix $Y$, then its cocycle would be a coboundary
$Y-\operatorname{Ad}r(g)Y$ in $\operatorname{Hom}(U,V)$.
This argument does not require $Y$ to be symplectic.
At $c$, normalization would give $2Y=0$.
Thus $Y=0$, which would imply $z=0$.
This proves nonsplitting and exactly the claimed
\texttt{BlockExt} condition.
\end{proof}

\begin{theorem}[RH Case II support for the rational points]
\label{thm:rh-four-rational-support}
For every original rational \texttt{KPair} satisfying
\texttt{Star}, in either coefficient characteristic,
the actual normalized pair on a fitted deep tower
lies in the support of the common ordinary Hecke
image ideal of Theorem~\ref{thm:rh-four-image-ideal}.
The tower CM field can be chosen unramified at any
prescribed finite set of rational primes outside the
ramification set $S$ of the original datum.
\end{theorem}

\begin{proof}
Choose the original labeling and initial tower by
Theorem~\ref{thm:rh-four-lat-shape} and its tame
tower lemma. All three fresh primes can avoid
the prescribed finite set. Perform the rational
lattice normalization, preserving the characteristic
polynomials and the exact labeled inertia characters.

Apply the deep residual-chain construction of
Section~\ref{sec:rh-case-II-chain}. Its auxiliary
primes, chosen by Chebotarev, can also avoid that
finite set. This fixes the finite starting tower
$T_1$ and its residual bridge before the numerical
degree threshold is chosen. Use the uniform
tower-bound theorem of Section~\ref{subsec:rh-four-uniform-tower-bounds}
to choose a cyclotomic layer with
$d\ge30(C+1)^2$. The same $C$ bounds all the
required cocycle spaces, including those after
the final coefficient enlargement. Only ramification
above $\ell$ is added by the cyclotomic layers.
The deep residual-chain theorem supplies the
arithmetic seed and the three links on this very
layer. Transport the actual normalized rational
pair by Lemma~\ref{lem:rh-four-restr-q}.

The common coefficient realization of
Section~\ref{subsec:rh-four-common-space} puts the
characteristic-zero seed, the actual target,
all residual links and the two ordinary links
over one coefficient DVR. Its finite-type space
gives the seed and target maps and the original
polynomial link maps. Form the common ideal $J$.
By the seed-local-zero assertion every prime below
the kernel of the characteristic-zero seed contains
$J$ on the coefficient ring. In particular the
residual seed has \texttt{SomeIn}; the component
propagation theorem gives \texttt{AllIn} there
and across all three polynomial links. It applies
in either characteristic and includes the possibly
nonsplit residue at the actual rational target.
Thus $J$ vanishes under the target coefficient
map. The restriction identities and the original
inertial labels are unchanged. This is the required
rational-point support conclusion, without asserting
the full \texttt{PMExAlt.aut} predicate.
\end{proof}

\begin{theorem}[The point-arithmetic inputs for the RH integral comparison]
\label{thm:rh-four-point-premises}
The preceding suppliers imply (F), (E), and (A) in
Section~\ref{sec:rh-symplectic-length}, for the actual
remaining RH Case II datum with $e\ge2$. In particular
the strict good-prime endpoint window suffices; no
additional assertion of strict windows at ramified
places or at the coefficient prime is used.
\end{theorem}

\begin{proof}
First every point used by the curve and endpoint criteria
is admitted by the original rational \texttt{KPair}.
In characteristic zero its flag and exact labels are
provided by Section~\ref{sec:rh-symplectic-point-rigidity};
in characteristic $\ell$ use the proper flag scheme
in the proof of Theorem~\ref{thm:rh-sp4-curve-criterion}.
The projective flag scheme has a point over a finite
extension of the fraction field of the curve DVR,
and the valuative criterion extends it integrally.
An integral symplectic basis adapted to its isotropic
flag exists by unimodularity of the symplectic form
and primitivity of the flag submodules (choose a basis
of the Lagrangian and its paired dual, correcting
the latter by a symmetric elementary transformation).
The graded characters have the prescribed inertial
values and residual values on the full decomposition
group: on the marked element $\mathrm{ph}$ this
is the specified specialization of (SR1), and
its unramified image generates the prime-to-$\ell$
part of the unramified quotient. On inertia the
labels are prescribed. A continuous character to a
finite residue field has order prime to $\ell$,
so the unramified prime-to-$\ell$ generator and
inertia determine its reduction. The small
coefficient and killing conditions come from the
actual universal model. Thus the earlier rational
lattice and support arguments apply to these
points, and Lemma~\ref{lem:rh-four-block-extension}
also applies to their block-diagonal endpoints.

\emph{Characteristic-prime exclusion (F).}
Suppose a forbidden characteristic-prime curve
of Theorem~\ref{thm:rh-sp4-curve-criterion} existed.
For an absolutely irreducible rank-four curve
take its rational representation. For a
block-diagonal curve with irreducible first
block replace it by the nonsplit extension of
Lemma~\ref{lem:rh-four-block-extension}, which
has exactly the same characteristic polynomials.
Use Theorem~\ref{thm:rh-four-rational-support}.
In the target coefficient map on the deep tower,
every ordinary inertial character is one: the
point has characteristic $\ell$ and the original
residual diagonal is trivial on the relative
inertia by \texttt{Ok.ineL}. The coefficient
uniformizer is zero as well. The finite-weight
clause of the common ideal now forces its
coefficient ring to evaluate through the finite
residue field, since the image is a domain.
Consequently the characteristic polynomial of
every element of $G_F$ equals that of the
prescribed residual diagonal there. For an
arbitrary $g\in G_{\mathbb Q}$ take a positive
power in $G_F$. All its powered eigenvalues
are algebraic over the finite residue field;
hence its eigenvalues themselves are algebraic
over it. The characteristic-polynomial
coefficients of the original point are integral
in the curve DVR. An element of that DVR
algebraic over the finite constant field equals
its residue lift; thus all original characteristic
polynomials, and in particular traces, equal
the prescribed residual ones on $G_{\mathbb Q}$.
The same conclusion holds for the block
diagonal point before forming the extension.
This contradicts the trace-algebra argument
(CF1)--(CF2) on the actual normalized curve,
which forbids both absolutely irreducible rank
four and an irreducible coordinate block when
traces are the residual diagonal traces. Both
forbidden types are excluded; the integral
curve criterion gives (F).

\emph{Characteristic-zero endpoint (E).}
If an original block-diagonal point had an
irreducible first block, again form the nonsplit
extension and apply rational-point support.
Its ordinary characters have the exact
prescribed arithmetic weights. The purity
part of Theorem~\ref{thm:rh-four-image-ideal}
gives equal complex absolute values for its
eigenvalues at good primes on the CM tower
field. At any rational prime outside the
finite ramification and exceptional sets take
a positive power of its arithmetic Frobenius
belonging to the CM field's absolute group;
it is a power of the corresponding unramified
Frobenius there. Equality of the absolute
values of the powered eigenvalues implies
equality before powering. This applies also
to the original block-diagonal representation
because the characteristic polynomials agree
at every group element. In particular at one
such prime $p>1$ all four absolute values
lie in a common strict half-power window
(indeed they are equal). But (CF3) gives
the ratio of the two absolute block
determinants exactly $p^2$. The strict
window would make that ratio strictly
less than $p^2$. This contradiction gives
(E).

\emph{Fixed-field automorphic attachment (A).}
Fix once and for all $K'$ from the four fields
of Theorem~\ref{thm:rh-symplectic-spectrum},
with its specified small split prime. Write
$q'$ for its odd ramified prime, which lies
outside $S$. For each irreducible original
point take the tower of
Theorem~\ref{thm:rh-four-rational-support},
choosing its CM field $L$ unramified at $q'$.
The field $K'$ is not changed as the point,
tower or coefficient extension changes.

We check the restriction irreducibility needed
for solvable descent, rather than assuming it.
If an absolutely irreducible representation
over an algebraically closed field of
characteristic zero becomes reducible on a
normal finite-index subgroup $H$ with abelian
quotient, it has a nontrivial self twist by a
character of the quotient. Indeed its restriction
to $H$ is semisimple: the sum of the finitely
many conjugates of a simple $H$-submodule
is semisimple and stable under the full
group, so is the full space. Then the
commutant of $H$ has dimension greater
than one. The quotient acts on it by
conjugation and splits it into character
eigenspaces. Its trivial eigenspace consists
of scalars by irreducibility. A nonzero
eigenvector of a nontrivial character is an
intertwiner with the associated self twist,
and is invertible by Schur's lemma.
Taking determinants shows that this
character has order dividing four.
Since $\ell$ is odd, reduction is injective
on these character values.

Apply this to $H=G_L$. Reduction of a
self-twist identity permutes the four
prescribed residual diagonal characters
(use the semisimplified residual
characteristic polynomials, equivalently
the Brauer--Nesbitt theorem). On $G_L$
these characters are pairwise distinct
by \texttt{Ok.g3} and the self-twist
character is trivial, so the permutation
must fix each character. The reduced
self-twist character, hence the character
itself, is trivial, a contradiction.
This proves irreducibility over $L$.

The same argument applies with $H=G_{LK'}$,
since $LK'/\mathbb Q$ is abelian. A self-twist
character trivial on this group is either
trivial on $G_L$, which was just excluded,
or restricts to the nontrivial quadratic
character of $LK'/L$. The latter is ramified
at $q'$ because $K'$ is ramified there and
$L$ is unramified there. On the other hand
the reduced character is a quotient of
two of the four original residual
characters, all unramified at $q'$.
Injectivity of reduction on the fourth
roots of unity rules out this ramification.
Thus restriction to $LK'$ is irreducible too.

The common ordinary Hecke ideal now supplies
cuspidal regular algebraic attachment over $L$
for $r|_{G_L}\otimes\kappa$, and untwisting
by the algebraic character gives the original
representation. In the source normalization
the polarization is the original one:
$\kappa\kappa^c\tau\varepsilon^3=1$,
so the twisted representation is
conjugate self dual with multiplier
$\varepsilon^{-3}$. Equivalently the
untwisted representation is polarized
with multiplier $\tau$; the symmetric
pairing is obtained from $J_{\rm sp}$ by
inserting $r(c)$ in its second argument,
since $\tau(c)=-1$ and $r(c)^2=1$.
This gives the polarized regular
algebraic cuspidal attachment required
for the classical soluble-base-change
and descent theorem. Apply Lemma 2.2.2
of Barnet-Lamb--Gee--Geraghty--Taylor,
\emph{Potential automorphy and change of weight},
Ann.\ Math.\ 179 (2014), 501--609, first to
$LK'/L$ and then to $LK'/K'$ (both soluble
Galois extensions of CM fields).
In each application the representation
over the larger field is irreducible,
as just proved. The theorem gives
polarized regular algebraic cuspidal
attachment over $K'$ to the original
restriction of $r$. These are ordinary
solvable base change and descent, not
a new residual automorphy lifting
assumption. We may use the same coefficient
isomorphism throughout. This is precisely
(A) of the integral comparison, with the
fixed quadratic field needed for the
polynomial trace-counting estimate.
\end{proof}

\paragraph{Consequence and scope.}
Substitution in Theorem~\ref{thm:rh-sp4-conditional-length}
gives, from the preceding suppliers, the individual remaining
RH Case II bound
\[
    (\log\ell)\ell(\ell,e,\psi)
                    \le (\ell N(e+1))^C .
\]
Here $C$ is the absolute exponent of the previously proved
integral comparison; the degrees and class-group bounds of
the auxiliary towers do not enter this numerical exponent.
This argument is restricted to the actual remaining RH
characters and is not the full original general Dirichlet
\texttt{PMExAlt.aut}. In particular it does not prove
\texttt{PotUnl} or the strict windows at all prescribed local
places. The good-prime window and fixed-field attachment
are the sufficient replacements used here. The weighted
product on the final common windows and the other independent
arithmetic and geometric suppliers remain separate.
